% !TeX program = xelatex
% 用 XeLaTeX 编译两遍，以更新目录与页码。
\documentclass[UTF8,a4paper,zihao=-4,fontset=fandol]{ctexart}
\usepackage[left=22mm,right=22mm,top=23mm,bottom=23mm,headheight=15pt,headsep=7mm,footskip=11mm]{geometry}
\usepackage{amsmath,amssymb}
\usepackage{array,booktabs,longtable}
\usepackage[table]{xcolor}
\usepackage{enumitem,needspace,fancyhdr,lastpage}
\usepackage[unicode,colorlinks,linkcolor=ink,urlcolor=sourceblue]{hyperref}
\usepackage{bookmark}

\definecolor{ink}{HTML}{1B3549}
\definecolor{sourceblue}{HTML}{285570}
\definecolor{quiet}{HTML}{53616A}
\definecolor{headerfill}{HTML}{EAF0F4}
\arrayrulecolor{black!37}

\hypersetup{
  pdftitle={哈工大模式识别与机器学习往年题总览},
  pdfsubject={2019 至 2025 年往年期末题与回忆题},
  pdfdisplaydoctitle=true,
  bookmarksopen=true
}
\setcounter{secnumdepth}{0}
\setcounter{tocdepth}{2}
\renewcommand{\contentsname}{目录}
\ctexset{
  section={format=\Large\bfseries\color{ink}, beforeskip=1.1em plus .2em, afterskip=.55em},
  subsection={format=\large\bfseries\color{ink}, beforeskip=1em plus .2em, afterskip=.35em},
  subsubsection={format=\normalsize\bfseries\color{ink}, beforeskip=.7em plus .15em, afterskip=.26em}
}
\setlength{\parindent}{0pt}
\setlength{\parskip}{.42em plus .08em minus .06em}
\setlength{\emergencystretch}{1.5em}
\clubpenalty=10000
\widowpenalty=10000
\raggedbottom
\setlist[itemize]{leftmargin=1.6em, topsep=.3em, itemsep=.42em, parsep=0pt}
\setlist[enumerate]{leftmargin=1.85em, label=\textbf{\arabic*.}, topsep=.3em, itemsep=.45em, parsep=0pt}
\setlength{\tabcolsep}{5pt}
\setlength{\LTleft}{0pt}
\setlength{\LTright}{0pt plus 1fill}
\setlength{\LTpre}{.35em}
\setlength{\LTpost}{.6em}
\newcolumntype{L}[1]{>{\raggedright\arraybackslash}p{\dimexpr#1\linewidth-2\tabcolsep\relax}}
\newcommand{\reporttitle}[1]{{\huge\bfseries\color{ink}#1\par}\vspace{.5em}}
\makeatletter
% Use URI actions for local paths with percent-encoded Unicode names.
\newcommand{\sourcelink}[2]{%
  \begingroup
  \let\%\@percentchar
  \edef\source@link{\noexpand\hyper@linkurl{\unexpanded{#2}}{#1}}%
  \source@link
  \endgroup
}
\makeatother

\pagestyle{fancy}
\fancyhf{}
\fancyhead[L]{\footnotesize\textcolor{quiet}{模式识别与机器学习}}
\fancyhead[R]{\footnotesize\textcolor{quiet}{往年题总览}}
\fancyfoot[C]{\footnotesize\textcolor{quiet}{第 \thepage\ 页\enspace/\enspace 共 \pageref*{LastPage}\ 页}}
\renewcommand{\headrulewidth}{.35pt}

\begin{document}

\reporttitle{哈工大模式识别与机器学习往年题总览}

本地资料共覆盖 \textbf{2019 秋至 2025 秋补／缓考的 9 套题}。其中 2020 春是四页 A 卷，其余为列题稿或回忆版。2022 春的 Word、Markdown 是同一套。2024 秋的 Word、PDF 互相补题。2024／2025 的简答参考答案可按卷对照。回忆中缺失的题干、选项和数据均在各题下标出。\par

\textbf{最贴近这次复习的两套是 2024 秋和 2025 秋补／缓考。} 两套都问到回归与 MSE、逻辑回归、SVM、RNN。2024 秋另有熵与 KL、普通自编码器、GMM／EM，2025 秋另有 KNN、DBSCAN、LSTM 和 Transformer。老师已明确说本次不考强化学习。本次模拟试题范本还明确出现了 VAE 采样生成、高斯朴素贝叶斯和 RNN 整句分类。 \sourcelink{file:///D:/Bhy_code/\%E4\%B8\%93\%E4\%B8\%9A\%E8\%AF\%BE/\%E5\%93\%88\%E5\%B0\%94\%E6\%BB\%A8\%E5\%B7\%A5\%E4\%B8\%9A\%E5\%A4\%A7\%E5\%AD\%A6\%E4\%BA\%BA\%E5\%B7\%A5\%E6\%99\%BA\%E8\%83\%BD\%E4\%B8\%93\%E4\%B8\%9A\%E8\%AF\%BE/\%E6\%A8\%A1\%E5\%BC\%8F\%E8\%AF\%86\%E5\%88\%AB\%E4\%B8\%8E\%E6\%9C\%BA\%E5\%99\%A8\%E5\%AD\%A6\%E4\%B9\%A0/PPT\%EF\%BC\%88\%E5\%88\%98\%E8\%BF\%9C\%E8\%B6\%85\%EF\%BC\%89/\%E6\%A8\%A1\%E6\%8B\%9F\%E8\%AF\%95\%E9\%A2\%98\%E8\%8C\%83\%E6\%9C\%AC-2026-11-3.pdf}{本次课程的模拟试题范本}\par

\vspace{.7em}
\tableofcontents
\clearpage
\section{跨卷反复出现的考点}

次数按下文的 \textbf{9 套本地试题}计。同一套多次提问算一次，同一套的多个版本也算一次。统计可定位的题目。只把明确出现的题意计入，对开放题按题干指定的模型归入相应主题。\par

\Needspace{5\baselineskip}
\begingroup
\small
\renewcommand{\arraystretch}{1.17}
\begin{longtable}{L{.19}L{.385}L{.425}}
\toprule
\rowcolor{headerfill}
\textbf{考点} & \textbf{出现套数与题号} & \textbf{可以练到什么程度} \\
\midrule
\endfirsthead
\toprule
\rowcolor{headerfill}
\textbf{考点} & \textbf{出现套数与题号} & \textbf{可以练到什么程度} \\
\midrule
\endhead
\midrule
\endfoot
\bottomrule
\endlastfoot
\textbf{聚类算法、GMM 与 EM} & \textbf{7 套}。2019 第 7 题，2021 秋第 4、6 题，2021 深度第 4 题，2022 春二（3），2022 秋第 3 题，2024 第三部分第 9 题，2025 简答第 2 题。 & 写 K-means 目标、分配与更新。说清 GMM 的 E 步算责任度、M 步更新权重、均值和协方差。2021 秋还要求比较 K-means 与 GMM，并处理协方差奇异。2025 要比较 K-means 与 DBSCAN。 \\
\addlinespace[.16em]
\textbf{贝叶斯分类} & \textbf{6 套}。2019 第 2、4 题，2021 秋第 3 题，2021 深度第 1 题，2022 春一（1），2022 秋第 5 题，2023 第 1 题。 & 掌握最小错误率决策、后验比较和朴素贝叶斯条件独立。2019 要证明最优性、写等协方差高斯的后验和边界。2022 秋再次问高斯情形何时产生线性决策面。 \\
\addlinespace[.16em]
\textbf{SVM 与核技巧} & \textbf{5 套}。2019 第 3、8 题，2021 秋第 5 题，2022 秋第 6 题，2024 单选第 2 题及第三部分第 4、5 题，2025 判断第 6 题及简答第 3 题。 & 会在图上指出超平面、支持向量和几何间隔。写出硬／软间隔的约束，解释对偶、松弛变量、惩罚系数和核函数分别解决什么问题。 \\
\addlinespace[.16em]
\textbf{逻辑回归} & \textbf{5 套}。2019 第 5 题，2021 秋第 7 题，2022 秋第 6 题，2024 第三部分第 3 题，2025 简答第 1、4 题。 & 会写 Sigmoid、正类概率和交叉熵。能证明常用阈值下的决策边界是线性的，并同线性回归、感知机、SVM 比较。 \\
\addlinespace[.16em]
\textbf{RNN、LSTM、GRU 与序列模型} & \textbf{5 套}。2020 第 3.4 题（选做的时序去噪），2021 深度第 9 题（PixelRNN 比较），2022 春第六题，2024 单选第 4 题及第三部分第 8 题，2025 判断第 2、4、7 题及简答第 6 题。 & 给输入和初始状态时能写出每步隐状态及输出。2024 特地指定\textbf{从右向左}，要看图确认是用前一个还是后一个状态。LSTM 复习门控和细胞状态，GRU 题练比较。 \\
\addlinespace[.16em]
\textbf{信息熵、信息增益和 KL} & \textbf{4 套}。2019 第 1 题，2021 秋第 2 题，2022 秋第 2 题，2024 第三部分第 6 题。 & 会算熵、条件熵、信息增益和指定方向的 KL。决策树题要用各子集的样本占比加权。2024 明确要求以 2 为底、保留两位小数。 \\
\addlinespace[.16em]
\textbf{Transformer 或 ViT} & \textbf{4 套}。2022 春五（2），2023 第 7 题，2024 判断第 3 题残片，2025 判断第 5 题及简答第 7 题。 & 2025 的具体问法是用文字讲自注意力、位置编码，以及为什么 RNN 通常无需额外位置编码。2022 春、2023 的问法更开放，涉及结构改进。 \\
\addlinespace[.16em]
\textbf{HMM} & \textbf{3 套}。2022 春一（2），2023 第 2 题，2024 单选第 4 题及判断第 2 题。 & 2022 春和 2023 都指向课件上的 HMM 例题，完整参数和问法存在缺口。2024 还问 HMM 与 RNN 的差别、HMM 能否预测。 \\
\addlinespace[.16em]
\textbf{PCA／FDA} & \textbf{3 套}。2019 第 6 题，2022 春二（2），2023 第 6 题。 & 2019 问 PCA 推导和两个应用，2022 春、2023 问 PCA 与 FDA 的区别。2021 深度第 3 题另有相关的内积矩阵与协方差矩阵谱关系。 \\
\end{longtable}
\endgroup

卷里也反复有\textbf{过拟合、数据划分和优化}。可以对照 2019 第 3 题、2021 秋第 1 题、2022 春第三至四题、2022 秋第 4 题、2023 第 4–5 题、2024 单选第 1 题和线性回归推导、2025 多项式回归简答。\textbf{最接近的计算题组合}是 2024 第三部分第 1–2 题与 2025 简答第 4、8 题，两套都要求把预测、误差、损失和更新式说清楚。\par

\Needspace{5\baselineskip}
\section{逐套逐题}

\Needspace{9\baselineskip}
\subsection{2019 年秋《机器学习》期末（9 题）}

现存材料逐题列了问法，共 9 道简答、推导或结合实验的分析题，未标分值。以下沿用原题号。\sourcelink{./\%E5\%93\%88\%E5\%B7\%A5\%E5\%A4\%A72019\%E7\%A7\%8B\%E6\%9C\%9F\%E6\%9C\%AB\%E8\%80\%83\%E8\%AF\%95.docx}{哈工大2019秋期末考试.docx}\par

\begin{enumerate}[start=1]
\item \textbf{信息增益与决策树边界。} 定义信息增益，用例子说明用它切分的好处。判断决策树的决策边界是否线性，并举例说明。
\item \textbf{最小错误率判别。} 给出使分类错误概率最小的最优判别方法，并证明最优性。
\item \textbf{过拟合。} 说明什么是过拟合，再分别解释决策树与 SVM 怎样降低过拟合风险。
\item \textbf{同协方差高斯分类。} 沿用第 2 题的判别准则。二元类别服从伯努利分布，两类的类条件分布为协方差相同的高斯分布。求决策边界方程和类别后验概率。
\item \textbf{逻辑回归与交叉熵。} 写出逻辑回归的优化目标，并指出把目标解释成交叉熵时，交叉熵中的两个概率分布分别是什么。
\item \textbf{PCA。} 推导 PCA，并给出两个应用例子。
\item \textbf{GMM、EM 与协方差奇异。} 写出 GMM 优化的目标和 EM 主要步骤。说明实验中协方差矩阵行列式为 0 的原因及解决方法。
\item \textbf{特征映射与线性可分性。} 解释低维空间线性不可分的样本为何可能在高维线性可分，举例并说明自己做过的实验是否用到这一思想。
\item \textbf{内积的应用。} 给出机器学习中三处使用内积的例子，解释内积在各自场景中的作用。
\end{enumerate}

\Needspace{9\baselineskip}
\subsection{2020 年春《模式识别与深度学习》试题 A（4 页）}

卷面共 100 分，考查两组工程问题。正文用第 1–5 题编号，展开为 10 个可独立定位的计分单元，包含选答项。题干、样例图、公式和卷积结构均按原 PDF 页核对。\sourcelink{./\%E5\%93\%88\%E5\%B7\%A5\%E5\%A4\%A72020\%E6\%98\%A5\%E6\%A8\%A1\%E5\%BC\%8F\%E8\%AF\%86\%E5\%88\%AB\%E4\%B8\%8E\%E6\%B7\%B1\%E5\%BA\%A6\%E5\%AD\%A6\%E4\%B9\%A0\%E8\%AF\%95\%E5\%8D\%B7\%E8\%AF\%95\%E9\%A2\%98.pdf}{哈工大2020春模式识别与深度学习试卷试题.pdf}\par

\Needspace{5\baselineskip}
\subsubsection{第一部分 小天体着陆区域选择（50 分）}

\Needspace{4\baselineskip}
\textbf{第 1 题 地形分类与岩石检测（第 1 页，必答 30 分）}\par

给出日本隼鸟 1 号拍摄的系川小行星灰度图。需依据岩石颗粒大小与散布规律识别三类地形，A 为平坦区，B 为碎石区，C 为大块岩石。约束是样本采集困难、探测器在轨算力有限。\par

\begin{itemize}
\item \textbf{1.1（1）特征计算。} 为灰度图 $I$ 中以 $(i,j)$ 为中心、宽 $w$ 高 $h$ 的矩形窗口设计特征映射，输出能近似表征岩石颗粒的向量 $x$。可组合多种特征，需说明计算思路及理由。
\item \textbf{1.1（2）地形分类。} 用 $x$ 设计一种或多种方法来区分三类地形，解释分类流程及选型理由。
\item \textbf{1.1（3）错分控制。} 结合地形\textup{III}图，说明如何避免把大块岩石表面误认成平坦地面，并设计验证实验。
\end{itemize}

\Needspace{4\baselineskip}
\textbf{第 2 题 小天体地形识别与着陆决策（第 2 页，三选一，20 分，多选无效）}\par

\begin{itemize}
\item \textbf{2.1 新颖性检测。} 设计方法发现区别于已知类型的新地形，再把新地形自动分成若干典型类别。交代算法思路和实验方法。
\item \textbf{2.2 阴影检测。} （1）考虑小行星自转产生的影子变化，说明如何识别和统计阴影与岩石的空间关系，并估计光照方向。（2）提出抑制阴影干扰、或反过来利用阴影变化改善地形分类的方法。
\item \textbf{2.3 着陆安全评分。} 对区域 $\Omega$ 设计可解释的评分 $V_\Omega=f(H,R,C,I,T,N,E)$。给定的七个输入依次为岩石尺寸分布直方图、由三维模型得到的表面粗糙度、五类地形标签 $C\in\{1,2,3,4,5\}$、受着陆时间与纬度影响的照明条件、反映表层和浅表层特性的热惯量、探测器动力学性能，以及着陆机构的地表适应能力。分别说明（1）模型形式、（2）参数估计、（3）验证方法。
\end{itemize}

\Needspace{5\baselineskip}
\subsubsection{第二部分 手机拍照图像增强（50 分）}

第 3.1、3.2 题为必答，共 30 分。第 3.3、3.4、4、5 题选答至 20 分。选答超过 20 分时，按原题顺序只计前 20 分。\par

\Needspace{4\baselineskip}
\textbf{第 3 题 手机拍照图像去噪（第 3 页）}\par

给出 LG G4 手机的有噪照片和对应清晰参考图。图 (c) 的网络是由多个 Conv 3×3 + ReLU 模块串联的去噪 CNN。\par

\begin{itemize}
\item \textbf{3.1 感受野与复杂度（必答 15 分）。} 在手机算力有限的条件下，求图示网络达到感受野 27 所需的层数。除控制层数外，提出至少三种平衡感受野与计算复杂性的做法，可调整网络结构，并结合去噪任务比较优缺点。
\item \textbf{3.2 损失函数与训练稳定性（必答 15 分）。} 训练数据为 $\{(y_i,x_i)\}_{i=1}^{N}$，去噪网络输出为 $\operatorname{CNN}(y_i)$。设计损失函数，解释其合理性。另提出改善视觉效果、保留丰富细节和提高训练稳定性的方案。
\item \textbf{3.3 训练性能（选答 10 分）。} 分别从网络设计、学习算法、集成学习等角度提出改进去噪网络训练性能的方法。
\item \textbf{3.4 视频去噪（选答 10 分）。} 在图 (c) 的去噪网络上加入时序神经网络模块。画出具体结构并解释设计依据。
\end{itemize}

\Needspace{4\baselineskip}
\textbf{第 4 题 手机图像超分辨率（第 3 页，选答 20 分）。} 给定不成对的低分辨率（LR）和高分辨率（HR）图像集合，用 GAN 学习图像超分辨率。说明深度卷积网络结构、损失函数和训练方法。\par

\Needspace{4\baselineskip}
\textbf{第 5 题 手机图像场景分类（第 3–4 页，选答 20 分）。} 为便于针对不同场景使用专门的图像增强模型，列出手机拍照会遇到的场景，说明分类标准及理由。分析它与一般图像分类相比有哪些特殊要求，以及应怎样调整网络结构和损失函数。\par

\Needspace{9\baselineskip}
\subsection{2021 年秋《机器学习》期末（8 题）}

现存材料列了 8 道简答或推导题，未标分值。按原题号整理如下。\sourcelink{./\%E5\%93\%88\%E5\%B7\%A5\%E5\%A4\%A72021\%E7\%A7\%8B\%E6\%9C\%BA\%E5\%99\%A8\%E5\%AD\%A6\%E4\%B9\%A0\%E6\%9C\%9F\%E6\%9C\%AB\%E8\%AF\%95\%E9\%A2\%98.docx}{哈工大2021秋机器学习期末试题.docx}\par

\begin{enumerate}[start=1]
\item \textbf{参数正则化与先验。} 解释参数正则化和参数先验的联系，以及在模型参数估计中加入正则化的目的。
\item \textbf{条件熵。} 写出条件熵定义，举出课程中一个使用它的例子。说明使用条件熵的好处，并给出直观解释。
\item \textbf{朴素贝叶斯。} 说明基本假设及它带来的好处。如果实际数据满足这一假设，判断分类器是否最优，并说明依据。
\item \textbf{K-means 与 GMM 的 EM。} 比较两种方法的共同点、差异和各自的优缺点。写出 GMM 的 E 步更新式和 M 步中需要最大化的期望目标。
\item \textbf{线性可分 SVM。} 写出硬间隔 SVM 的目标函数、约束和解的形式。说明为何求解对偶问题及其好处。
\item \textbf{GMM 协方差矩阵不可逆。} 解释 EM 优化中出现不可逆协方差的原因，并说明遇到这种情况时会采用什么处理方法。
\item \textbf{逻辑回归的决策边界。} 判断逻辑回归分类器是否为线性分类器，并给出证明或理由。
\item \textbf{内积的应用。} 举出三个课程中使用向量或函数内积的场景，分别解释它的作用。
\end{enumerate}

\Needspace{9\baselineskip}
\subsection{2021 年《深度学习与模式识别》考试（回忆 10 题）}

资料未注明学期或分值。第 1、2、8、9、10 题部分细节未记录，以下注明各自缺项。\sourcelink{./\%E5\%93\%88\%E5\%B7\%A5\%E5\%A4\%A72021\%E6\%B7\%B1\%E5\%BA\%A6\%E5\%AD\%A6\%E4\%B9\%A0\%E4\%B8\%8E\%E6\%A8\%A1\%E5\%BC\%8F\%E8\%AF\%86\%E5\%88\%AB\%E8\%80\%83\%E8\%AF\%95\%E9\%A2\%98.docx}{哈工大2021深度学习与模式识别考试题.docx}\par

\begin{enumerate}[start=1]
\item \textbf{癌症案例中的贝叶斯判别。} 回忆指向 PPT 中判断有无癌症的原例题，未留下疾病先验、检测结果、相关概率或具体问法。
\item \textbf{极大似然估计。} 回忆只说明题目与 PPT 例题类似，分布、样本和待估参数均未记录。
\item \textbf{内积矩阵与协方差矩阵。} 写出两个矩阵的表达式，比较相同与不同之处，判断特征值是否相同并证明。
\item \textbf{谱聚类。} 解释谱聚类的基本思想、“谱”的含义，以及为什么最后还要用 K-means 聚类。
\item \textbf{Sigmoid 与 ReLU。} 说明为什么在相关神经网络中后来常以 ReLU 替代 Sigmoid。
\item \textbf{Adam 的矩估计。} 考查 Adam 的一阶矩和二阶矩。是否要求写更新式或偏差校正，回忆未说明。
\item \textbf{GAN 的模型与应用。} 选一种 GAN 模型，说明其具体用法和应用方向。
\item \textbf{卷积神经网络。} 回忆称此题与 2020 年的卷积题类似，未记录输入、网络结构或子问。相关原题为 2020 年春 A 卷第 3 题。
\item \textbf{BERT 与 PixelRNN／PixelCNN。} 比较两者的异同，至少涉及训练方式和上下文定义。原回忆另提到一个比较维度，但未记录内容。
\item \textbf{WGAN 的上确界约束。} 回忆提到上确界下方有“范数小于 1”的约束，问模型如何满足这一约束。约束对象和空间名称只能按回忆理解，资料未留下原式。
\end{enumerate}

\Needspace{9\baselineskip}
\subsection{2022 年春　模式识别与深度学习}

来源 \sourcelink{./\%E5\%93\%88\%E5\%B7\%A5\%E5\%A4\%A72022\%E6\%98\%A5\%E6\%A8\%A1\%E5\%BC\%8F\%E8\%AF\%86\%E5\%88\%AB\%E4\%B8\%8E\%E6\%B7\%B1\%E5\%BA\%A6\%E5\%AD\%A6\%E4\%B9\%A0\%E6\%9C\%9F\%E6\%9C\%AB\%E8\%AF\%95\%E9\%A2\%98.docx}{哈工大2022春模式识别与深度学习期末试题.docx}、\sourcelink{./\%E5\%93\%88\%E5\%B7\%A5\%E5\%A4\%A72022\%E6\%98\%A5\%E6\%A8\%A1\%E5\%BC\%8F\%E8\%AF\%86\%E5\%88\%AB\%E4\%B8\%8E\%E6\%B7\%B1\%E5\%BA\%A6\%E5\%AD\%A6\%E4\%B9\%A0\%E6\%9C\%9F\%E6\%9C\%AB\%E8\%AF\%95\%E9\%A2\%98.md}{哈工大2022春模式识别与深度学习期末试题.md}。两版题序、题意和已记录的参数一致，合并为一套。共 8 大题、17 小问。回忆稿注明每小问 4–7 分，多为 5 或 6 分，未逐题注明分值。\par

\Needspace{5\baselineskip}
\begingroup
\small
\renewcommand{\arraystretch}{1.17}
\begin{longtable}{L{.14}L{.86}}
\toprule
\rowcolor{headerfill}
\textbf{原题号} & \textbf{题目概要} \\
\midrule
\endfirsthead
\toprule
\rowcolor{headerfill}
\textbf{原题号} & \textbf{题目概要} \\
\midrule
\endhead
\midrule
\endfoot
\bottomrule
\endlastfoot
一（1） & 给定类别先验概率和类条件概率，求最小错误率贝叶斯分类器。回忆稿未保留概率数值。 \\
\addlinespace[.16em]
一（2） & HMM 课件中“朋友在外地”的例题变式。转移矩阵第二行改为 $(0.5, 0.5)$。回忆稿未保留其余参数、完整观测序列或具体所求。 \\
\addlinespace[.16em]
一（3） & 从递归贝叶斯学习的角度概述卡尔曼滤波的原理。 \\
\addlinespace[.16em]
二（1） & 解释课件“拉普拉斯性质 1”的作用和各符号的含义，并证明该性质。回忆稿未抄出性质公式。 \\
\addlinespace[.16em]
二（2） & 从相似图出发，比较 PCA 与基于 FDA（LDA）的方法的特点。 \\
\addlinespace[.16em]
二（3） & 概述谱聚类的思路，说明最后一步采用 K-means 的原因。 \\
\addlinespace[.16em]
三（1） & 以中国人、欧洲人和非洲人的人脸识别为背景。方案 A 按这三组人脸依次作训练、验证、测试集。方案 B 先按老人、青年、小孩、男人、女人分组，再按 50\%、30\%、20\% 划分训练、验证、测试集。指出两种方案的问题并给出自己的划分建议。 \\
\addlinespace[.16em]
三（2） & 根据学习曲线提出下一步优化措施并解释理由。回忆者记为高方差情形，并提到第二章课件有相应的学习曲线。文件未附题图。 \\
\addlinespace[.16em]
四（1） & 解释 MLP 难以学习深层网络的原因，给出 3 种使感知器从浅层扩展到深层的改进方法。 \\
\addlinespace[.16em]
四（2） & 介绍至少 4 种使用动量的随机优化方法，并比较其区别。 \\
\addlinespace[.16em]
五（1） & 给出 4 种在参数量不显著增加的情况下扩大 CNN 感受野的方法。 \\
\addlinespace[.16em]
五（2） & 讨论 Transformer 的发展能为 CNN 的改进提供哪些借鉴。 \\
\addlinespace[.16em]
六 & 比较 LSTM 和 GRU 的门控单元。 \\
\addlinespace[.16em]
七（1） & 说明 Improved WGAN 如何满足 Lipschitz 约束，比较它与 WGAN 的约束方式。 \\
\addlinespace[.16em]
七（2） & 列举 4 种把人脸图像逆映射到预训练人脸生成模型（如 StyleGAN）隐空间的方法，并介绍一种应用。 \\
\addlinespace[.16em]
八（1） & 以语义分割为例，从标注成本、标注样本数等方面比较小样本学习、弱监督学习和自监督学习。回忆稿未记清大题标题。 \\
\addlinespace[.16em]
八（2） & 在图像分类网络的基础上，说明如何调整网络结构与损失函数来实现端到端目标检测，并举例。 \\
\end{longtable}
\endgroup

\Needspace{9\baselineskip}
\subsection{2022 年秋　机器学习}

原稿列出 7 题。见 \sourcelink{./\%E5\%93\%88\%E5\%B7\%A5\%E5\%A4\%A7\%E6\%9C\%AC\%E9\%83\%A82022\%E7\%A7\%8B\%E6\%9C\%BA\%E5\%99\%A8\%E5\%AD\%A6\%E4\%B9\%A0\%E6\%9C\%9F\%E6\%9C\%AB\%E8\%AF\%95\%E9\%A2\%98.docx}{哈工大本部2022秋机器学习期末试题.docx}。\par

\Needspace{5\baselineskip}
\begingroup
\small
\renewcommand{\arraystretch}{1.17}
\begin{longtable}{L{.14}L{.86}}
\toprule
\rowcolor{headerfill}
\textbf{原题号} & \textbf{题目概要} \\
\midrule
\endfirsthead
\toprule
\rowcolor{headerfill}
\textbf{原题号} & \textbf{题目概要} \\
\midrule
\endhead
\midrule
\endfoot
\bottomrule
\endlastfoot
1 & 给出有监督学习和无监督学习的定义，各列举 3 个算法。 \\
\addlinespace[.16em]
2 & 解释熵、条件熵和互信息。根据某决策树节点上候选属性 $X_1,X_2$ 各自的分类结果，选择更好的切分属性并说明理由。原稿提到分类结果“如下”，实际未附计数或图表。 \\
\addlinespace[.16em]
3 & 写出 K-means 的步骤与优化目标，并说明关键步骤如何使目标得到优化。 \\
\addlinespace[.16em]
4 & 解释过拟合，结合自己做过的实验说明，并给出 2 种应对方法。 \\
\addlinespace[.16em]
5 & 两类样本的类条件分布均为高斯，指出贝叶斯分类器的决策面成为线性决策面所需的分布条件。 \\
\addlinespace[.16em]
6 & 从损失函数、决策平面等方面比较 SVM 与逻辑回归的异同。 \\
\addlinespace[.16em]
7 & 谈机器学习课程的学习心得。 \\
\end{longtable}
\endgroup

\Needspace{9\baselineskip}
\subsection{2023 年　模式识别与深度学习}

原稿列出 12 题，未标春秋学期。见 \sourcelink{./\%E5\%93\%88\%E5\%B7\%A5\%E5\%A4\%A72023\%E6\%A8\%A1\%E5\%BC\%8F\%E8\%AF\%86\%E5\%88\%AB\%E4\%B8\%8E\%E6\%B7\%B1\%E5\%BA\%A6\%E5\%AD\%A6\%E4\%B9\%A0\%E6\%9C\%9F\%E6\%9C\%AB.docx}{哈工大2023模式识别与深度学习期末.docx}。\par

\Needspace{5\baselineskip}
\begingroup
\small
\renewcommand{\arraystretch}{1.17}
\begin{longtable}{L{.14}L{.86}}
\toprule
\rowcolor{headerfill}
\textbf{原题号} & \textbf{题目概要} \\
\midrule
\endfirsthead
\toprule
\rowcolor{headerfill}
\textbf{原题号} & \textbf{题目概要} \\
\midrule
\endhead
\midrule
\endfoot
\bottomrule
\endlastfoot
1 & 最小错误率贝叶斯决策。原稿未保留给定条件、数值或具体所求。 \\
\addlinespace[.16em]
2 & HMM 课件中的天气与行动例题。原稿未保留状态转移、发射与初始概率或小问。 \\
\addlinespace[.16em]
3 & 解释卡尔曼滤波各个参数的含义。原稿未列出具体符号。 \\
\addlinespace[.16em]
4 & 写出 ReLU 激活函数及其 3 种变体的数学表达式。 \\
\addlinespace[.16em]
5 & 优化算法。原稿只有这个主题，未注明具体算法或小问。 \\
\addlinespace[.16em]
6 & 比较 FDA 与 PCA。回忆者提到可能要求从相似图解释，该条件未获确认。 \\
\addlinespace[.16em]
7 & 说明如何改进 Transformer（ViT）。原稿未给定具体改进目标或模型设置。 \\
\addlinespace[.16em]
8 & 从重建精度与推理速度比较 StyleGAN、pSp 和 PTI。 \\
\addlinespace[.16em]
9 & 写出 Faster R-CNN 与第二种目标检测算法的损失函数，分析主要区别。第二种算法原文写作 “DERT”，疑指 DETR。 \\
\addlinespace[.16em]
10 & 从标注成本和样本数量两方面比较全监督、半监督与小样本学习。 \\
\addlinespace[.16em]
11 & 设计一个基于模式识别的监考系统，至少提出 4 项功能，并给出各功能的具体算法或实现。 \\
\addlinespace[.16em]
12 & 介绍自己实现过的一种模式识别系统，说明关键问题及解决思路。回忆稿未记清其余要求。 \\
\end{longtable}
\endgroup

\Needspace{9\baselineskip}
\subsection{2024 秋《模式识别与机器学习 A》期末考试}

卷面合计 60 分。单选 5 道，每题 1 分，判断 10 道，每题 1 分，计算分析 9 道，每题 5 分。以下以回忆 DOCX 的题号为准。PDF 第 1 页的 8 道可辨认主观题与 DOCX 重合，排列顺序不同。PDF 未忆及的第九道主观题可由 DOCX 的自编码器题补足。参考答案复现了 DOCX 的 9 道主观题，硬间隔、软间隔和 RNN 三张配图与回忆 DOCX 一致。\sourcelink{./2024\%E5\%B9\%B4\%E7\%A7\%8B\%E6\%A8\%A1\%E5\%BC\%8F\%E8\%AF\%86\%E5\%88\%AB\%E4\%B8\%8E\%E6\%9C\%BA\%E5\%99\%A8\%E5\%AD\%A6\%E4\%B9\%A0A\%E6\%9C\%9F\%E6\%9C\%AB\%E8\%80\%83\%E8\%AF\%95\%EF\%BC\%88\%E5\%9B\%9E\%E5\%BF\%86\%E7\%89\%88\%EF\%BC\%89.docx}{2024年秋模式识别与机器学习A期末考试（回忆版）.docx} · \sourcelink{./2024\%E5\%B9\%B4\%E7\%A7\%8B\%E6\%A8\%A1\%E5\%BC\%8F\%E8\%AF\%86\%E5\%88\%AB\%E4\%B8\%8E\%E6\%9C\%BA\%E5\%99\%A8\%E5\%AD\%A6\%E4\%B9\%A0A\%E5\%9B\%9E\%E5\%BF\%86.pdf}{2024年秋模式识别与机器学习A回忆.pdf} · \sourcelink{./\%E5\%93\%88\%E5\%B7\%A5\%E5\%A4\%A72024\%E5\%B9\%B4\%E7\%A7\%8B\%2B2025\%E5\%B9\%B4\%E7\%A7\%8B\%E7\%BC\%93\%E8\%80\%83\%E7\%AE\%80\%E7\%AD\%94\%E9\%A2\%98_\%E5\%8F\%82\%E8\%80\%83\%E7\%AD\%94\%E6\%A1\%88.docx}{哈工大2024年秋+2025年秋缓考简答题\_参考答案.docx}\par

\Needspace{5\baselineskip}
\subsubsection{单选题}

\begin{itemize}
\item \textbf{第 1 题（DOCX）} 识别哪项措施不适合处理过拟合。记得有“减少训练数据量”这一选项，其余选项未载。
\item \textbf{第 2 题（DOCX）} 根据“把低维特征映射到高维特征空间，以找到分离超平面”的描述识别模型，涉及 SVM 的特征映射。
\item \textbf{第 3 题} DOCX 未记下题干或选项。
\item \textbf{第 4 题（DOCX）} 辨别 HMM 与 RNN 的概念和差别。残存线索涉及 HMM 是否基于神经网络、含隐变量的时间序列建模和 RNN 参数更新，具体选项未载。
\item \textbf{第 5 题（DOCX）} 将若干模型划分为有监督学习和无监督学习，具体模型名单未载。
\item \textbf{无编号线索（PDF 第 1 页）} 作者以“把输出层的值转成概率分布用什么函数（Softmax）”举例说明选择题难度。材料未给该例题号，无法确认其是否对应 DOCX 缺失的第 3 题。
\end{itemize}

\Needspace{5\baselineskip}
\subsubsection{判断题}

\begin{itemize}
\item \textbf{第 1 题（DOCX）} 判断验证集是否在训练时用于更新模型参数，以及它是否用于报告模型结果。
\item \textbf{第 2 题（DOCX）} 判断 HMM 不能用于数据序列预测的说法。
\item \textbf{第 3 题（DOCX，仅残片）} 涉及 Transformer 自注意力的实现，作者另记有“长短效”字样，完整命题与具体对象均未载。
\item \textbf{第 4–10 题} DOCX 未记下题干。PDF 第 1 页以“学习率过低会导致梯度更新跨度过大”为错误命题举例，未给题号或完整题面。
\end{itemize}

\Needspace{9\baselineskip}
\subsubsection{计算分析题}

\Needspace{5\baselineskip}
\begingroup
\small
\renewcommand{\arraystretch}{1.17}
\begin{longtable}{L{.12}L{.67}L{.21}}
\toprule
\rowcolor{headerfill}
\textbf{DOCX 题号} & \textbf{回忆题意和数据} & \textbf{PDF 第 1 页对应项} \\
\midrule
\endfirsthead
\toprule
\rowcolor{headerfill}
\textbf{DOCX 题号} & \textbf{回忆题意和数据} & \textbf{PDF 第 1 页对应项} \\
\midrule
\endhead
\midrule
\endfoot
\bottomrule
\endlastfoot
1 & \textbf{线性回归与梯度下降。} 已知目标值 $y$ 和输入向量 $x$，自行定义权重、偏置、学习率等，写出误差平方和，推导损失梯度和参数迭代更新式。 & 第 1 项 \\
\addlinespace[.16em]
2 & \textbf{多项式回归与 MSE。} 说明多项式回归基本原理。模型为 $\hat y=0.8x^2$，测试集第 6、7 个样本分别为 $(x,y)=(4.0,14.8),(5.0,22.0)$。求测试 MSE，采用 PDF 明确记下的 $1/N$ 系数。训练集在表中以省略号表示。 & 第 2 项，未保留测试数据 \\
\addlinespace[.16em]
3 & \textbf{逻辑回归。} 写出 Sigmoid 的表达式，比较逻辑回归和线性回归。题干说明 Sigmoid 把线性得分映射到 0–1，用作正类概率。 & 第 6 项 \\
\addlinespace[.16em]
4 & \textbf{硬间隔 SVM。} 结合二维灰白两类样本的可分散点图，解释分离超平面、支持向量和样本点到分离超平面的几何间隔。图中绿色实线标为 $H_1$，两侧虚线表示间隔边界。 & 第 7 项，几何距离的对象记忆较模糊 \\
\addlinespace[.16em]
5 & \textbf{软间隔 SVM 与核技巧。} 结合样本进入间隔或落在线的另一侧的配图，说明训练数据无法完全线性分开时软间隔最大化的约束和基本原理。另说明核技巧如何处理非线性分类。 & 第 8 项 \\
\addlinespace[.16em]
6 & \textbf{信息熵与 KL 散度。} 给出 $P=(0.1,0.3,0.6)$、$Q=(0.1,0.2,0.7)$，分别求 $H(P)$、$H(Q)$ 和 $D_{\mathrm{KL}}(P\Vert Q)$，以 2 为底，保留两位小数。DOCX 的“相对熵”方向由 PDF 第 3 项“P 对 Q”补明。 & 第 3 项 \\
\addlinespace[.16em]
7 & \textbf{自编码器。} 编码器把输入向量 $x$ 转为 $z$，解码器由 $z$ 产生重构 $x'$。比较 $x'$ 和 $x$ 的长度并说明原因，再用图像重构解释工作原理和用途。 & PDF 未记下该题的内容 \\
\addlinespace[.16em]
8 & \textbf{反向时序的 RNN 前向计算。} 根据示意图写出当前时刻输出 $o_t$ 的表达式，可自行定义参数。横向箭头是从右向左，$h_t$ 接收 $x_t$ 与 $h_{t+1}$ 的信息，而各输出节点由同一时刻的隐状态指向。 & 第 4 项，方向由 DOCX 题干和配图确认 \\
\addlinespace[.16em]
9 & \textbf{GMM 与 EM。} 说明 EM 的核心思想和具体步骤，指出 GMM 聚类中计算哪些量对应 E 步，更新哪些量对应 M 步。 & 第 5 项 \\
\end{longtable}
\endgroup

PDF 第 1 页还记载了当年有班级的考前讲座，以及回忆者称该卷未考 PCA。第 2 页的逻辑回归、MLP、CNN、CLIP/CoOp、StyleGAN 和 LoRA 属于课程实验回顾。这些是当年回忆者的补充记录。\sourcelink{./2024\%E5\%B9\%B4\%E7\%A7\%8B\%E6\%A8\%A1\%E5\%BC\%8F\%E8\%AF\%86\%E5\%88\%AB\%E4\%B8\%8E\%E6\%9C\%BA\%E5\%99\%A8\%E5\%AD\%A6\%E4\%B9\%A0A\%E5\%9B\%9E\%E5\%BF\%86.pdf}{2024年秋模式识别与机器学习A回忆.pdf}\par

\Needspace{9\baselineskip}
\subsection{2025 秋《模式识别与机器学习》补／缓考}

回忆 PDF 共 3 页。补／缓考卷面 100 分，单选 5 道共 10 分、判断 10 道共 10 分、简答 8 道共 80 分。PDF 作者另记正常期末卷面 80 分，并认为题目范围与难度接近。该评价来自回忆材料。参考答案的后半部分复现以下可辨认的 7 道简答，仍缺第 5 题。\sourcelink{./ML25\%E7\%A7\%8B\%E8\%A1\%A5\%E7\%BC\%93\%E8\%80\%83\%E5\%9B\%9E\%E5\%BF\%86\%E7\%89\%88.pdf}{ML25秋补缓考回忆版.pdf} · \sourcelink{./\%E5\%93\%88\%E5\%B7\%A5\%E5\%A4\%A72024\%E5\%B9\%B4\%E7\%A7\%8B\%2B2025\%E5\%B9\%B4\%E7\%A7\%8B\%E7\%BC\%93\%E8\%80\%83\%E7\%AE\%80\%E7\%AD\%94\%E9\%A2\%98_\%E5\%8F\%82\%E8\%80\%83\%E7\%AD\%94\%E6\%A1\%88.docx}{哈工大2024年秋+2025年秋缓考简答题\_参考答案.docx}\par

\Needspace{5\baselineskip}
\subsubsection{单选题（PDF 第 1 页）}

\begin{itemize}
\item \textbf{第 1 题（仅选项）} KNN 的 $k$ 越小，模型越复杂，也越容易过拟合。回忆版只保留这一选项。
\item \textbf{第 2 题（仅主题和选项）} 比较 K-means 与 DBSCAN。记住的选项声称 K-means 对初始聚类中心不敏感，不同初始化得到相同结果，回忆者标为错误。
\item \textbf{第 3 题} 一个特征的均值为 1000，其余特征均值为 1，问归一化的主要目的。回忆记录的答案指向避免单个特征主导权重更新。
\item \textbf{第 4–5 题} 题干和选项均未载。
\end{itemize}

\Needspace{5\baselineskip}
\subsubsection{判断题（PDF 第 1–2 页）}

\begin{itemize}
\item \textbf{第 1 题} 判断感知机能否拟合非线性问题。
\item \textbf{第 2 题} 判断 LSTM 是否通过 Sigmoid 把门控值限制在 0–1，进而控制各门的信息保留强度。
\item \textbf{第 3 题} 判断 KNN 的 $k$ 越大是否会使偏差越小、方差越大。
\item \textbf{第 4 题} 判断 RNN 能处理任意长度序列是否主要源于各时间步共享参数，回忆文字使用“拟合”。
\item \textbf{第 5 题} 判断 Transformer 是否主要依靠递归和卷积实现。
\item \textbf{第 6 题} 判断 SVM 能否通过核技巧把非线性分类转化为线性分类问题。
\item \textbf{第 7 题} 判断 LSTM 能否利用输入门、遗忘门、输出门应对普通 RNN 的长期信息遗忘。
\item \textbf{第 8–10 题} 第 2 页第 8 项仅写其余题记不清，未载第 8–10 题的题干。
\end{itemize}

\Needspace{9\baselineskip}
\subsubsection{简答题（PDF 第 2–3 页）}

\Needspace{5\baselineskip}
\begingroup
\small
\renewcommand{\arraystretch}{1.17}
\begin{longtable}{L{.21}L{.79}}
\toprule
\rowcolor{headerfill}
\textbf{原题号与页码} & \textbf{小问及可辨认数据} \\
\midrule
\endfirsthead
\toprule
\rowcolor{headerfill}
\textbf{原题号与页码} & \textbf{小问及可辨认数据} \\
\midrule
\endhead
\midrule
\endfoot
\bottomrule
\endlastfoot
1，第 2 页 & \textbf{感知机与逻辑回归。} （1）比较两者模型形式及分类判定方式。（2）问哪一种能找到最大几何间隔的分离超平面。 \\
\addlinespace[.16em]
2，第 2 页 & \textbf{K-means 与 DBSCAN。} （1）说明 K-means 的主要流程与缺点。（2）客户按收入、消费水平拟分为高收入高消费、高收入低消费、低收入高消费、低收入低消费四类，数据中含噪声。选择 K-means 或 DBSCAN 并说明原因，再解释 DBSCAN 如何识别与处理噪声。 \\
\addlinespace[.16em]
3，第 2 页 & \textbf{SVM。} （1）解释支持向量。（2）对不能完全线性分开的情形，选择硬间隔或软间隔，说明如何实现。 \\
\addlinespace[.16em]
4，第 2 页 & \textbf{损失函数。} （1）由预测 $\hat y$ 和真值 $y$ 求 MSE，原始数据未载。（2）说出逻辑回归训练时采用的损失函数。 \\
\addlinespace[.16em]
5，第 2 页 & \textbf{缺失。} PDF 只留了两处小问占位，没有可辨认题干。 \\
\addlinespace[.16em]
6，第 2 页 & \textbf{RNN 与 LSTM。} （1）初始状态记为 $h_0$，输入序列 $x=(2,-1)$，按 $h_t=\tanh(0.8x_t+0.1h_{t-1})$ 写出 $h_1,h_2$ 的表达式。回忆 PDF 未写 $h_0$ 的具体数值，式中的输入按时刻取 $x_t$。（2）说明 LSTM 主要解决普通 RNN 的什么问题，以及门控和记忆如何起作用。 \\
\addlinespace[.16em]
7，第 2–3 页 & \textbf{Transformer。} （1）用文字说明自注意力怎么实现，题目明确无需写公式。（2）位置编码起什么作用。（3）RNN 是否需要位置编码，说明理由。 \\
\addlinespace[.16em]
8，第 3 页 & \textbf{多项式回归。} （1）写出 MSE 损失。（2）求以 MSE 为损失时对权重 $\omega$ 的梯度 $\partial J/\partial \omega$。（3）解释高阶多项式在训练集拟合好、测试集效果差的现象及成因。 \\
\end{longtable}
\endgroup

参考答案第 4（1）问中的真值 $(1,2,3)$ 和预测 $(1,2.5,2)$ 是解题示例。第 6（1）问后的 $h_0=0$ 是代算时的假设。两组数字均不能补作原题条件。\sourcelink{./\%E5\%93\%88\%E5\%B7\%A5\%E5\%A4\%A72024\%E5\%B9\%B4\%E7\%A7\%8B\%2B2025\%E5\%B9\%B4\%E7\%A7\%8B\%E7\%BC\%93\%E8\%80\%83\%E7\%AE\%80\%E7\%AD\%94\%E9\%A2\%98_\%E5\%8F\%82\%E8\%80\%83\%E7\%AD\%94\%E6\%A1\%88.docx}{哈工大2024年秋+2025年秋缓考简答题\_参考答案.docx}\par

\clearpage
\section{对照答案时容易混的几处}

\begin{itemize}
\item \textbf{2024 多项式回归的原数据可直接复算。} $x=4,5$ 时预测分别为 $12.8,20.0$，真值是 $14.8,22.0$，所以两个误差都是 $2$，MSE 为 $4$。2025 简答第 4（1）问的原始数据没有留存。参考答案中的三元向量是另举的例子。
\item \textbf{2024 第 6 题要保留 KL 的方向和对数底。} 对给定的 $P,Q$，以 2 为底算得 $H(P)\approx1.30$、$H(Q)\approx1.16$、$D_{\mathrm{KL}}(P\Vert Q)\approx0.04$，单位为 bit。
\item \textbf{感知机、逻辑回归和 SVM 解决的问题不同。} 2025 简答第 1（2）问涉及最大几何间隔。按课程里的常规定义，显式最大化几何间隔的是硬间隔 SVM。软间隔通过松弛变量处理违反间隔的样本，核技巧通过替换内积来使用非线性特征。在线性原空间中单独加软间隔，决策函数仍是线性的。软间隔支持向量也可能落在间隔内部，因而不能把它们一律写成 $y_i(w^\top x_i+b)=1$。
\item \textbf{RNN 的方向要随题看。} 2024 图题从右向左，因此 $h_t$ 依赖 $h_{t+1}$。2025 给出正向递推，保留未知的 $h_0$ 时，$h_1=\tanh(1.6+0.1h_0)$，$h_2=\tanh(-0.8+0.1h_1)$。原题初值的数值没有留在回忆版中。
\item \textbf{自注意力和顺序。} 对无位置编码、无掩码的全局自注意力，输入顺序改变时，输出跟着作同样的重排。这个性质叫置换等变。写 2025 位置编码题时，表述成「位置编码给模型提供顺序信息」即可。
\item \textbf{把普通 AE、VAE 和 GAN 分开。} 2024 第三部分第 7 题是普通自编码器的重构与维度，早年的 GAN 问法另见 2020 选做题、2021 深度、2022 春和 2023。VAE 的生成在本次范本里有选择题。 \sourcelink{file:///D:/Bhy_code/\%E4\%B8\%93\%E4\%B8\%9A\%E8\%AF\%BE/\%E5\%93\%88\%E5\%B0\%94\%E6\%BB\%A8\%E5\%B7\%A5\%E4\%B8\%9A\%E5\%A4\%A7\%E5\%AD\%A6\%E4\%BA\%BA\%E5\%B7\%A5\%E6\%99\%BA\%E8\%83\%BD\%E4\%B8\%93\%E4\%B8\%9A\%E8\%AF\%BE/\%E6\%A8\%A1\%E5\%BC\%8F\%E8\%AF\%86\%E5\%88\%AB\%E4\%B8\%8E\%E6\%9C\%BA\%E5\%99\%A8\%E5\%AD\%A6\%E4\%B9\%A0/PPT\%EF\%BC\%88\%E5\%88\%98\%E8\%BF\%9C\%E8\%B6\%85\%EF\%BC\%89/\%E6\%A8\%A1\%E6\%8B\%9F\%E8\%AF\%95\%E9\%A2\%98\%E8\%8C\%83\%E6\%9C\%AC-2026-11-3.pdf}{本次课程的模拟题范本}
\end{itemize}

\Needspace{5\baselineskip}
\section{参考链接与资料缺口}

\begin{itemize}
\item \textbf{2023 年 A 课程回忆帖。} \sourcelink{./\%E5\%8F\%82\%E8\%80\%83\%E9\%93\%BE\%E6\%8E\%A5.md}{参考链接} 指向的知乎文章返回 403，正文未核到。它的标题写「模式识别与机器学习 A」，本地 2023 Word 写「模式识别与深度学习」，无法确认是否同一场。本地 Word 的 12 题已逐项列出。 \href{https://zhuanlan.zhihu.com/p/681977670}{知乎原链接}
\item \textbf{2024 年 A 课程回忆帖。} CSDN 原文链接未能打开。本地 PDF 第 2 页标明来自该文章的修改版本，PDF 里的题目线索已与同场 Word 回忆版合并。 \sourcelink{./2024\%E5\%B9\%B4\%E7\%A7\%8B\%E6\%A8\%A1\%E5\%BC\%8F\%E8\%AF\%86\%E5\%88\%AB\%E4\%B8\%8E\%E6\%9C\%BA\%E5\%99\%A8\%E5\%AD\%A6\%E4\%B9\%A0A\%E5\%9B\%9E\%E5\%BF\%86.pdf}{本地 PDF} · \href{https://blog.csdn.net/awawawawawawawaa/article/details/144856589?spm=1001.2014.3001.5502}{CSDN 原链接}
\item \textbf{2024 和 2025 回忆稿的缺号。} 2024 有线索的题号共 16 个，另有 PDF 记下的两条无题号示例。2025 秋补／缓考有线索的题号共 17 个，简答第 5 题仅剩两个占位。题目缺失不计入跨卷次数。
\item \textbf{Word 锁文件。} 目录里的文件名以 \texttt{\textasciitilde{}\$} 开头的文件是 Word 临时锁，共 162 字节。
\end{itemize}

\end{document}
